\NSPage{115}%
For the fixed background and initialization, derivatives through order
\NSi{NS-F-P115-001}{m} are bounded by
\NSi{NS-F-P115-002}{C_mq^{-K_m}(1+|\log q|)^{P_m}}. Finally, \eqref{eq:9.18} is
\eqref{eq:5.33} for the fixed differential polynomial
\NSFormula{NS-F-P115-003}{display}{none}
\[
\mathcal F(u,p)=\mathscr R(u,p)=\partial_tu+(u\mathbin{\cdot}\nabla)u-\Delta u+\nabla p,
\]
with \NSi{NS-F-P115-004}{\rho_j=h\sigma_j\to\infty}. The estimate \eqref{eq:9.18}
concerns the actual pointwise residual, after evaluating the auxiliary variables at
\NSi{NS-F-P115-005}{Y(r,t)}. The finite-stage flat term is uniform over its infinitely many
dyadic labels. The finite initialization cutoff only adds errors supported away from
\NSi{NS-F-P115-006}{q=0}, which also meet that fixed-stage flatness condition. These residual
estimates are uniform on the whole local domain: choose a fixed bounded interval in
\NSi{NS-F-P115-007}{X} containing all slow and annular correction supports. The estimates above
and \eqref{eq:5.41} apply there. Outside it, every finite state agrees with the exact exterior heat
field of \eqref{eq:4.29}, with its centrifugal pressure, so its residual is identically zero.

\noindent\textbf{Step 2: The summed fields and their Cartesian regularity.} Lemma~\ref{lem:5.4}
therefore gives potentials whose locally finite sum defines a velocity satisfying
\NSi{NS-F-P115-008}{\operatorname{div}u_{\mathrm{loc}}=0}, and
\NSFormula{NS-F-P115-009}{display}{9.20}
\begin{equation}
\forall m,N\qquad |\mathscr R(u_{\mathrm{loc}},p_{\mathrm{loc}})|_m=O(q^N)
\qquad(q\downarrow0).
\tag{9.20}\label{eq:9.20}
\end{equation}
The tail estimate compares the constructed field with one fixed finite state and its flat remainder.
The flat residuals separated at each finite stage are estimated through this comparison and are not
summed. To identify the potentials and pressure of the summed fields, let
\NSi{NS-F-P115-010}{\mathcal A_0,\mathcal B_0e_\theta,p_0} be the representatives of
\NSi{NS-F-P115-011}{U_0}: the realized slow base from \eqref{eq:5.27}, together with the finite
initialization and its pressure. Define
\NSFormula{NS-F-P115-012}{display}{9.21}
\begin{equation}
\begin{aligned}
\mathcal A&=\mathcal A_0+\sum_{j\geq1}\chi(a_jq)A_j,\\
\mathcal Be_\theta&=\mathcal B_0e_\theta+\sum_{j\geq1}\chi(a_jq)B_j,\\
p_{\mathrm{loc}}&=p_0+\sum_{j\geq1}\chi(a_jq)p_j,
&u_{\mathrm{loc}}&=\operatorname{curl}\mathcal A+\mathcal Be_\theta.
\end{aligned}
\tag{9.21}\label{eq:9.21}
\end{equation}
This gives the required vector potential, azimuthal field, and pressure. At the axis the base Stokes
streamfunctions have the form \NSi{NS-F-P115-013}{r^2} times a smooth function of
\NSi{NS-F-P115-014}{(r^2,z,t)}, so the base vector potential is smooth in Cartesian coordinates.
The base angular field is \NSi{NS-F-P115-015}{r} times such a scalar times
\NSi{NS-F-P115-016}{e_\theta}. All annular representatives vanish in a neighborhood of the axis
at each \NSi{NS-F-P115-017}{t<1}. Thus \eqref{eq:9.21} proves
Theorem~\ref{thm:3.1}\textup{(i)}.

\noindent\textbf{Step 3: One-sided regularity away from the singular point.} It remains to verify
the endpoint regularity away from \NSi{NS-F-P115-018}{q=0}. Fix a compact spatial set and
\NSi{NS-F-P115-019}{0<c<c'<q_*}. Only finitely many terms in the expansion in powers of
\NSi{NS-F-P115-020}{q^{2h}} and correction stages have nonzero cutoff factors on a neighborhood
of \NSi{NS-F-P115-021}{c\leq q\leq c'}, because both cutoff sequences tend to infinity. There are
also finitely many dyadic bands and slow labels on this compact set. The base profiles and their
Stokes streamfunctions have bounds for derivatives of every fixed order on the closed parameter
range \NSi{NS-F-P115-022}{-1\leq\eta\leq1}, by Theorem~\ref{thm:4.6} and
Proposition~\ref{prop:5.5}. The geometry chooses its labels, representatives, rounded frequencies,
and enlarged local torus rectangles from the closed \NSi{NS-F-P115-023}{\tau\geq0} range and
keeps these discrete choices fixed as \NSi{NS-F-P115-024}{\tau\to0}. The pulse inverse acts along
a path with the physical slow variables fixed, within \NSi{NS-F-P115-025}{\tau\geq0}.
Differentiating its fixed linear ODE gives the bounds at every derivative order. These bounds, the
quotient bound \eqref{eq:7.33} using the fixed positive primary amplitudes, and the radial and
temporal mean maps give finite one-sided bounds for every fixed derivative of each wave potential,
mean
\NSPage{116}%
vector potential, azimuthal mean component, and pressure whose cutoff is nonzero on the chosen
neighborhood. Their smooth zero extensions cover every support boundary.

Coordinate conversion has denominator \NSi{NS-F-P116-001}{L\geq1-2h>0}, so on
\NSi{NS-F-P116-002}{q\geq c} it preserves these bounds for derivatives of each fixed order. On
annular supports \NSi{NS-F-P116-003}{r} is bounded away from zero; at the axis use the preceding
Cartesian factors. The finite sum therefore has uniformly bounded Cartesian space--time derivatives
of every order up to \NSi{NS-F-P116-004}{\tau=0}, also inside the active shell. Bounding one
additional time derivative and applying the fundamental theorem of calculus gives the uniform
one-sided limit of each derivative. The same theorem on compact spatial subsets proves compatibility
with spatial derivatives; compatibility between limits of consecutive time derivatives follows by
integration in time. This proves Theorem~\ref{thm:3.1}\textup{(ii)} for
\NSi{NS-F-P116-005}{\mathcal A,\mathcal Be_\theta,p_{\mathrm{loc}}} themselves.

\noindent\textbf{Step 4: The exterior heat field and flat residual.} All annular wave and mean
correction potentials, azimuthal mean components, and pressures vanish beyond
\NSi{NS-F-P116-006}{X_b}. The zero total axial moments of the leading profile and of each
coefficient in the expansion in powers of \NSi{NS-F-P116-007}{q^{2h}} also make their Stokes
streamfunctions vanish in the common exterior. Consequently, for a fixed
\NSi{NS-F-P116-008}{X_{\mathrm{ext}}} beyond all slow supports,
\NSi{NS-F-P116-009}{\mathcal A=0} and \NSi{NS-F-P116-010}{\mathcal B=K}. The leading-profile pressure,
normalized to vanish at radial infinity, and the compactly supported pressure corrections of positive
order give precisely \NSi{NS-F-P116-011}{-\int_r^\infty K(\rho,\tau)^2\,d\rho/\rho} there. Using
\eqref{eq:4.29}, define
\NSFormula{NS-F-P116-012}{display}{none}
\[
\mathcal H_{\mathrm{ext}}(s)=2^Ac_\infty\mathcal H(4s),
\qquad K(r,\tau)=r^{-1-2h}\mathcal H_{\mathrm{ext}}(\tau/r^2).
\]
This is \eqref{eq:3.5}. The integral formula \eqref{eq:A.34} gives, for every integer
\NSi{NS-F-P116-013}{m\geq0},
\NSFormula{NS-F-P116-014}{display}{none}
\[
\sup_{s\geq0}|\mathcal H_{\mathrm{ext}}^{(m)}(s)|
\leq2^Ac_\infty4^m(h)_m(1+h)_m<\infty,
\]
where \NSi{NS-F-P116-015}{(b)_m} is the rising factorial and
\NSi{NS-F-P116-016}{(b)_0=1}. Indeed, the factor
\NSi{NS-F-P116-017}{(1+Zv)^{-h-m}} in that integral is at most one for every
\NSi{NS-F-P116-018}{Z,v\geq0}. Thus the required derivative bounds hold on the full nonnegative
half-line. By Lemma~\ref{lem:A.6}, \NSi{NS-F-P116-019}{K} solves the radial heat equation; its
centrifugal pressure makes the Navier--Stokes residual identically zero in the exterior. Together
with the background residual estimates on compact sets of profile variables in \eqref{eq:5.41} and
\eqref{eq:9.20}, this proves all of Theorem~\ref{thm:3.1}\textup{(iii)}.

\noindent\textbf{Step 5: The inner velocity growth.} Finally choose any fixed
\NSi{NS-F-P116-020}{X_{\mathrm{in}}\in(0,X_a)} where the positive leading azimuthal velocity is
\NSi{NS-F-P116-021}{e_0=E_0(X_{\mathrm{in}},0)>0}. All annular corrections vanish there. The
realized expansion in powers of \NSi{NS-F-P116-022}{q^{2h}} in \eqref{eq:5.1} has its leading
value plus a remainder \NSi{NS-F-P116-023}{O(q^{2h})} after removing the velocity factor
\NSi{NS-F-P116-024}{q^{-A}}, at the fixed point
\NSi{NS-F-P116-025}{(X,\eta)=(X_{\mathrm{in}},0)} in profile coordinates. At
\NSi{NS-F-P116-026}{z=0} one has \NSi{NS-F-P116-027}{\eta=0} and
\NSi{NS-F-P116-028}{q=\tau}, hence
\NSFormula{NS-F-P116-029}{display}{none}
\[
u_{\theta,\mathrm{loc}}(\sqrt{2X_{\mathrm{in}}\tau},0,0,1-\tau)
=\tau^{-A}\bigl(e_0+O(\tau^{2h})\bigr).
\]
This is Theorem~\ref{thm:3.1}\textup{(iv)}, and completes its proof.
\end{proof}

\section{Compact forcing and whole-space breakdown}\label{sec:10}

In this section we turn the fields supplied by Theorem~\ref{thm:3.1} into the whole-space solution
and force in Theorem~\ref{thm:1.1}. We write the local fields as
\NSFormula{NS-F-P116-030}{display}{none}
\[
u^{\mathrm{loc}}=\operatorname{curl}\mathcal A+\mathcal Be_\theta,
\qquad p^{\mathrm{loc}}.
\]
Here \NSi{NS-F-P116-031}{\mathcal A} and \NSi{NS-F-P116-032}{\mathcal B} are the representatives
constructed in \eqref{eq:9.21}, and \NSi{NS-F-P116-033}{\mathcal B} is independent of
\NSi{NS-F-P116-034}{\theta}. We first work at viscosity one. The local velocity has the growth
\eqref{eq:3.6} required for Theorem~\ref{thm:1.1};
\NSPage{117}%
we must give it zero initial datum and compact spatial support, while making its momentum residual
the restriction of a force in
\NSi{NS-F-P117-001}{C_c^\infty(\mathbb R^3\times(0,\infty);\mathbb R^3)}
(Proposition~\ref{prop:10.1} and Lemma~\ref{lem:10.3}).

We multiply the vector potential, azimuthal coefficient, and pressure by cutoffs equal to one near
\NSi{NS-F-P117-002}{(0,1)}. We then prove that every derivative of the resulting residual has a
uniform limit as \NSi{NS-F-P117-003}{t\uparrow1} (Lemma~\ref{lem:10.2}), and construct a force
attaining those limits from \NSi{NS-F-P117-004}{t>1} (Lemma~\ref{lem:10.3}). The equation gives
a uniform energy bound before time one (Lemma~\ref{lem:10.4}). Finally, comparison with any smooth
solution of bounded kinetic energy (Lemma~\ref{lem:10.5}) transfers the local velocity growth to that
solution, giving the contradiction needed for Theorem~\ref{thm:1.1}. Rescaling gives arbitrary
positive viscosity and the periodic construction (Corollary~\ref{cor:10.6}).

\subsection{Localization within the local domain.}\label{sec:10.1}
In this subsection we localize the velocity and pressure inside the domain of
Theorem~\ref{thm:3.1}. We first recall why the chosen vector potential
\NSi{NS-F-P117-005}{\mathcal A} is smooth at the axis (Theorem~\ref{thm:3.1}\textup{(i)}) and
vanishes in the heat exterior \eqref{eq:3.5}. These properties will allow us to take its curl after
localization and to estimate the force at time one (Lemma~\ref{lem:10.2}). For the background
coefficients \NSi{NS-F-P117-006}{U_n} in \eqref{eq:5.1}, the Stokes streamfunctions and their sum
are
\NSFormula{NS-F-P117-007}{display}{10.1}
\begin{equation}
\begin{gathered}
\mathcal S_n=q^{1-A+2nh}\int_0^XU_n(X',\eta)\,dX',
\qquad \mathcal S=\mathcal S_0+\sum_{n\geq1}\chi(c_nq)\mathcal S_n,\\
\mathcal A_{\mathrm{base}}=\frac{\mathcal S}r e_\theta.
\end{gathered}
\tag{10.1}\label{eq:10.1}
\end{equation}
The cutoffs \NSi{NS-F-P117-008}{\chi(c_nq)} are those used in
Proposition~\ref{prop:5.5}. By \eqref{eq:4.28} and \eqref{eq:5.10},
\NSi{NS-F-P117-009}{\int_0^\infty U_n(X,\eta)\,dX=0} for every
\NSi{NS-F-P117-010}{n\geq0}. Thus each radial integral in \eqref{eq:10.1} vanishes beyond the
common support of the axial coefficients. Near the axis,
\NSi{NS-F-P117-011}{\mathcal S=r^2a(r^2,z,t)} for a smooth scalar \NSi{NS-F-P117-012}{a}, and hence
\NSFormula{NS-F-P117-013}{display}{none}
\[
(\mathcal S/r)e_\theta=a(r^2,z,t)(-x_2,x_1,0).
\]
This formula gives the smooth Cartesian representative of the background potential.

Every wave velocity is the curl of its supported physical potential. For a prescribed axial mean
increment \NSi{NS-F-P117-014}{\gamma_{d,\mathrm{phys}}}, the mean construction uses the azimuthal
potential coefficient
\NSFormula{NS-F-P117-015}{display}{none}
\[
\Psi=r^{-1}\mathcal I_c(r\gamma_{d,\mathrm{phys}}),
\qquad \mathcal I_cg=\mathcal Ig-\chi_m\mathcal Jg,
\]
with the radial operations defined in \eqref{eq:8.4}. Beyond the source and the transition of
\NSi{NS-F-P117-016}{\chi_m}, we have \NSi{NS-F-P117-017}{\chi_m=1} and
\NSi{NS-F-P117-018}{\mathcal Ig=\mathcal Jg}, so \NSi{NS-F-P117-019}{\Psi=0}. These wave and
mean potentials, with their summation cutoffs, are the terms added to
\NSi{NS-F-P117-020}{\mathcal A_{\mathrm{base}}} in \eqref{eq:9.21}. The direct azimuthal mean
increments are included in \NSi{NS-F-P117-021}{\mathcal B}. Evaluation at
\NSi{NS-F-P117-022}{Y(r,t)} preserves their independence of \NSi{NS-F-P117-023}{\theta}, and
all the annular corrections vanish near the axis for each fixed \NSi{NS-F-P117-024}{t<1}. In
particular, for the fixed \NSi{NS-F-P117-025}{X_{\mathrm{ext}}} in
Theorem~\ref{thm:3.1}, the chosen potential satisfies
\NSi{NS-F-P117-026}{\mathcal A=0} on \NSi{NS-F-P117-027}{X\geq X_{\mathrm{ext}}}.

We now choose cutoffs supported inside the local domain
\NSi{NS-F-P117-028}{\Omega_* = \{\tau>0,\ q<q_*\}}, where
\NSi{NS-F-P117-029}{\tau=1-t}. Applying the spatial cutoff before taking the curl preserves
incompressibility, and the temporal cutoff gives zero initial datum.

\begin{proposition}\label{prop:10.1}
There are a compact set \NSi{NS-F-P117-030}{\mathcal K\subset\mathbb R^3}, a smooth vector field
\NSi{NS-F-P117-031}{u} and a smooth scalar field \NSi{NS-F-P117-032}{p} on
\NSi{NS-F-P117-033}{\mathbb R^3\times[0,1)}, with
\NSi{NS-F-P117-034}{\supp u(\mathord\cdot,t)\cup\supp p(\mathord\cdot,t)\subset\mathcal K} for every
\NSi{NS-F-P117-035}{0\leq t<1}, satisfying
\NSFormula{NS-F-P117-036}{display}{10.2}
\begin{equation}
\operatorname{div}u=0,
\qquad u(\mathord\cdot,t)=p(\mathord\cdot,t)=0
\qquad\text{for all sufficiently small }t\geq0.
\tag{10.2}\label{eq:10.2}
\end{equation}
They agree with \NSi{NS-F-P117-037}{u^{\mathrm{loc}},p^{\mathrm{loc}}} on a spatial neighborhood
of the origin for all times sufficiently close to \NSi{NS-F-P117-038}{1}.
\end{proposition}
\NSPage{118}%
\begin{proof}
To keep a fixed spatial cutoff inside the local domain, we need an upper bound for
\NSi{NS-F-P118-001}{q} in terms of \NSi{NS-F-P118-002}{z} and
\NSi{NS-F-P118-003}{\tau}. The coordinates in \eqref{eq:3.2} give such a bound uniformly in
\NSi{NS-F-P118-004}{r}. If \NSi{NS-F-P118-005}{1-\eta^2\geq1/2}, then
\NSi{NS-F-P118-006}{q\leq2\tau}; otherwise
\NSi{NS-F-P118-007}{|\eta|>2^{-1/2}} and
\NSi{NS-F-P118-008}{q\leq(\sqrt2|z|)^{1/D}}. Consequently
\NSFormula{NS-F-P118-009}{display}{10.3}
\begin{equation}
q\leq C_0\bigl(\tau+|z|^{1/D}\bigr).
\tag{10.3}\label{eq:10.3}
\end{equation}
We choose \NSi{NS-F-P118-010}{0<\tau_0<1/2} and \NSi{NS-F-P118-011}{z_0>0} with
\NSi{NS-F-P118-012}{C_0(\tau_0+z_0^{1/D})<q_*/2}. For any fixed
\NSi{NS-F-P118-013}{r_0>0}, we choose a smooth axisymmetric cutoff
\NSi{NS-F-P118-014}{0\leq\chi_x\leq1}, smooth as a function of
\NSi{NS-F-P118-015}{r^2,z}, such that
\NSFormula{NS-F-P118-016}{display}{none}
\[
\chi_x=1\quad(r\leq r_0/2,\ |z|\leq z_0/2),
\qquad \supp\chi_x\Subset\{r<r_0,\ |z|<z_0\}.
\]
We also choose a smooth function \NSi{NS-F-P118-017}{0\leq\chi_t\leq1} of
\NSi{NS-F-P118-018}{t} with
\NSFormula{NS-F-P118-019}{display}{none}
\[
\chi_t(t)=0\quad(1-t\geq\tau_0),
\qquad \chi_t(t)=1\quad(0\leq1-t\leq\tau_0/2),
\qquad c(x,t)=\chi_x(x)\chi_t(t).
\]
By \eqref{eq:10.3}, the support of \NSi{NS-F-P118-020}{c} for
\NSi{NS-F-P118-021}{t<1} lies in \NSi{NS-F-P118-022}{q<q_*/2}. In particular it stays a fixed
distance in the \NSi{NS-F-P118-023}{q} coordinate from the outer boundary of
\NSi{NS-F-P118-024}{\Omega_*}. On that domain we define
\NSFormula{NS-F-P118-025}{display}{10.4}
\begin{equation}
u=\operatorname{curl}(c\mathcal A)+c\mathcal Be_\theta,
\qquad p=cp^{\mathrm{loc}},
\tag{10.4}\label{eq:10.4}
\end{equation}
and extend by zero outside the cutoff support. The product rule gives
\NSFormula{NS-F-P118-026}{display}{none}
\[
u=cu^{\mathrm{loc}}+\nabla c\times\mathcal A,
\qquad \operatorname{div}u=\operatorname{div}\operatorname{curl}(c\mathcal A)
 +r^{-1}\partial_\theta(c\mathcal B)=0.
\]
The Cartesian representatives recalled above make the fields smooth at the axis. The cutoffs and
the margin from \NSi{NS-F-P118-027}{q=q_*} give smooth zero extension elsewhere. The support is
contained in \NSi{NS-F-P118-028}{\mathcal K=\supp\chi_x}, and the temporal cutoff gives \eqref{eq:10.2}.
Both cutoffs are one near \NSi{NS-F-P118-029}{(0,1)}, so the local fields are unchanged there.
\end{proof}

For \NSi{NS-F-P118-030}{0\leq t<1} we define
\NSFormula{NS-F-P118-031}{display}{10.5}
\begin{equation}
f=\mathscr R(u,p)=\partial_tu+(u\mathbin{\cdot}\nabla)u-\Delta u+\nabla p.
\tag{10.5}\label{eq:10.5}
\end{equation}
This definition includes every derivative of the cutoffs. Taking its divergence gives
\NSFormula{NS-F-P118-032}{display}{none}
\[
-\Delta p=\sum_{i,j=1}^3\partial_i\partial_j(u_iu_j)-\operatorname{div}f.
\]
Thus the localized pressure satisfies the Poisson equation with the corresponding force-divergence
term; the force may have nonzero divergence.

\subsection{Limits of force derivatives at the terminal time.}\label{sec:10.2}
The force in \eqref{eq:10.5} is so far defined only for \NSi{NS-F-P118-033}{t<1}. A smooth
extension requires compatible limits of all its derivatives as
\NSi{NS-F-P118-034}{t\uparrow1}. Flatness of the local residual in \eqref{eq:3.4} controls these
limits at the origin; away from the origin, the endpoint bounds in
Theorem~\ref{thm:3.1}\textup{(ii)} and the exact exterior heat flow \eqref{eq:3.5} control the terms
introduced by localization. In this subsection we first establish the derivative limits
(Lemma~\ref{lem:10.2}), then construct an extension with compact time support
(Lemma~\ref{lem:10.3}).

\begin{lemma}\label{lem:10.2}
For every spatial multi-index \NSi{NS-F-P118-035}{\alpha} and integer
\NSi{NS-F-P118-036}{j\geq0}, \NSi{NS-F-P118-037}{\partial_x^\alpha\partial_t^jf} converges
uniformly on \NSi{NS-F-P118-038}{\mathbb R^3} as
\NSi{NS-F-P118-039}{t\uparrow1}. There are functions
\NSi{NS-F-P118-040}{F_j\in C_c^\infty(\mathbb R^3;\mathbb R^3)}, all supported in
\NSi{NS-F-P118-041}{\mathcal K}, such that
\NSFormula{NS-F-P118-042}{display}{10.6}
\begin{equation}
\lim_{t\uparrow1}\partial_x^\alpha\partial_t^jf(x,t)=\partial_x^\alpha F_j(x),
\qquad \partial_x^\alpha F_j(0)=0.
\tag{10.6}\label{eq:10.6}
\end{equation}
These are the derivative limits of \NSi{NS-F-P118-043}{f} from
\NSi{NS-F-P118-044}{t<1}, compatible under spatial and time differentiation.
\end{lemma}
\NSPage{119}%
\begin{proof}
We first establish endpoint bounds away from the origin. On compact spatial sets with
\NSi{NS-F-P119-001}{q} bounded below and bounded away from \NSi{NS-F-P119-002}{q_*},
Theorem~\ref{thm:3.1}\textup{(ii)} bounds every physical derivative of the actual potentials and
the scalar fields. More precisely, on such a compact set \NSi{NS-F-P119-003}{K_0}, for
\NSi{NS-F-P119-004}{G=\mathcal A,\mathcal Be_\theta,p^{\mathrm{loc}}} and
\NSi{NS-F-P119-005}{t<t'<1}, the fundamental theorem of calculus gives
\NSFormula{NS-F-P119-006}{display}{none}
\[
\sup_{K_0}|\partial_x^\alpha\partial_t^jG(\mathord\cdot,t')
 -\partial_x^\alpha\partial_t^jG(\mathord\cdot,t)|
\leq|t'-t|\sup_{K_0\times[t,t']}|\partial_x^\alpha\partial_t^{j+1}G|.
\]
The last supremum is bounded independently of \NSi{NS-F-P119-007}{t,t'} near one. Thus each fixed
derivative is uniformly Cauchy as \NSi{NS-F-P119-008}{t\uparrow1}. All constants here may depend
on the positive lower bound for \NSi{NS-F-P119-009}{q}.

The argument using finitely many nonzero cutoff terms in Proposition~\ref{prop:9.9} supplies these
bounds on the closed parameter range \NSi{NS-F-P119-010}{-1\leq\eta\leq1}, for the actual
representatives and the smooth Cartesian factors in their near-axis representations, with the labels
and integer frequencies fixed. Multiplication by the fixed cutoffs and the finite differential
expression \eqref{eq:10.5} preserve the bounds.

The remaining region has \NSi{NS-F-P119-011}{r} bounded below and
\NSi{NS-F-P119-012}{q} sufficiently small, so \NSi{NS-F-P119-013}{X\geq X_{\mathrm{ext}}}
and the velocity is the exterior azimuthal field
\NSi{NS-F-P119-014}{K(r,\tau)e_\theta} of \eqref{eq:3.5}. In terms of the heat profile
\NSi{NS-F-P119-015}{\mathcal H} defined by \eqref{eq:A.32},
\NSFormula{NS-F-P119-016}{display}{10.7}
\begin{equation}
K(r,\tau)=c_\infty s^{-A}\mathcal H(2\tau/s),
\qquad s=r^2/2.
\tag{10.7}\label{eq:10.7}
\end{equation}
Thus the function in Theorem~\ref{thm:3.1} is
\NSi{NS-F-P119-017}{\mathcal H_{\mathrm{ext}}(Z)=2^Ac_\infty\mathcal H(4Z)}. For each
\NSi{NS-F-P119-018}{j}, differentiating the integral in \eqref{eq:A.32} is justified uniformly for
\NSi{NS-F-P119-019}{Z\geq0} by a constant times
\NSi{NS-F-P119-020}{e^{-v}v^{h+j}}. It follows that
\NSFormula{NS-F-P119-021}{display}{10.8}
\begin{equation}
|\partial_\tau^jK(r,\tau)|\leq C_jr^{-1-2h-2j},
\qquad
\left|\partial_\tau^j\frac{K(\rho,\tau)^2}{\rho}\right|
\leq C_j\rho^{-3-4h-2j}.
\tag{10.8}\label{eq:10.8}
\end{equation}
The pressure is normalized to vanish at radial infinity. On a compact positive-radius interval
\NSi{NS-F-P119-022}{r\in[r_-,r_+]}, with \NSi{NS-F-P119-023}{r_->0}, the second majorant is
integrable for \NSi{NS-F-P119-024}{\rho\geq r_-}. Dominated convergence therefore justifies
\NSFormula{NS-F-P119-025}{display}{none}
\[
\partial_\tau^jp^{\mathrm{loc}}(r,\tau)
=-\int_r^\infty\partial_\tau^j\left(\frac{K(\rho,\tau)^2}{\rho}\right)\,d\rho.
\]
It also gives a uniform bound and a one-sided limit for each derivative on compact positive-radius
intervals. Here \NSi{NS-F-P119-026}{\partial_t=-\partial_\tau}. Further radial derivatives of the
pressure follow by differentiating the lower endpoint; those of \NSi{NS-F-P119-027}{K} follow
directly from \eqref{eq:10.7}. This proves the required smooth one-sided extension to
\NSi{NS-F-P119-028}{\tau=0} using the integral on \NSi{NS-F-P119-029}{\tau\geq0}. In this
region \NSi{NS-F-P119-030}{\mathcal A=0}, so the localized fields involve only these smooth
exterior quantities and fixed cutoffs. The heat equation for \NSi{NS-F-P119-031}{K} and
\NSi{NS-F-P119-032}{\partial_rp^{\mathrm{loc}}=K^2/r} also show directly that the uncut exterior
residual is exactly zero.

These two regions cover the terminal slice away from the origin: if
\NSi{NS-F-P119-033}{z\neq0}, then \NSi{NS-F-P119-034}{q\geq|z|^{1/D}>0}; if
\NSi{NS-F-P119-035}{z=0} and \NSi{NS-F-P119-036}{x\neq0}, then
\NSi{NS-F-P119-037}{r>0} and the preceding exterior argument applies on a small neighborhood.
Hence every Cartesian space--time derivative of the force converges uniformly on compact spatial
sets avoiding the origin.

Near the origin the cutoffs equal one. We use \eqref{eq:3.4} on
\NSi{NS-F-P119-038}{X\leq X_{\mathrm{ext}}}, and the exact zero residual for larger
\NSi{NS-F-P119-039}{X}. For every \NSi{NS-F-P119-040}{\alpha,j,N} this gives, on a fixed
neighborhood of \NSi{NS-F-P119-041}{(0,1)},
\NSFormula{NS-F-P119-042}{display}{10.9}
\begin{equation}
|\partial_x^\alpha\partial_t^jf(x,t)|
\leq C_{\alpha,j,N}\bigl(\tau+|z|^{1/D}\bigr)^N.
\tag{10.9}\label{eq:10.9}
\end{equation}
For a prescribed tolerance, we first choose a small ball and then a late enough time so that this
bound is below the tolerance there. A finite cover of the rest of \NSi{NS-F-P119-043}{\mathcal K} gives
uniform convergence outside that ball. All these derivatives vanish off
\NSi{NS-F-P119-044}{\mathcal K}. They are therefore uniformly Cauchy on
\NSi{NS-F-P119-045}{\mathbb R^3}, and their limits vanish at the origin.
\NSPage{120}%
We define \NSi{NS-F-P120-001}{F_j} to be the uniform limit for
\NSi{NS-F-P120-002}{\alpha=0}. Passing to the limit in the fundamental theorem of calculus along
spatial coordinate segments identifies every other limit as
\NSi{NS-F-P120-003}{\partial_x^\alpha F_j}. Thus \NSi{NS-F-P120-004}{F_j} is smooth and has
the stated support. Passing to the limit in the corresponding identity in time gives
\NSFormula{NS-F-P120-005}{display}{10.10}
\begin{equation}
\partial_x^\alpha\partial_t^jf(x,t)
=\partial_x^\alpha F_j(x)-\int_t^1\partial_x^\alpha\partial_t^{j+1}f(x,s)\,ds.
\tag{10.10}\label{eq:10.10}
\end{equation}
This proves compatibility of the derivative limits under spatial and time differentiation and
completes \eqref{eq:10.6}.
\end{proof}

The limits \NSi{NS-F-P120-006}{F_j} give the Taylor coefficients needed to continue the force
through time one. We realize them with time cutoffs whose supports shrink with the derivative order,
so that the extension is also compactly supported in time. Such a smooth compactly supported force
satisfies the decay conditions in \cite[(5)]{ref-13}; the zero datum satisfies its initial-data
conditions as well.

\begin{lemma}\label{lem:10.3}
The force in \eqref{eq:10.5} extends to
\NSi{NS-F-P120-007}{f\in C_c^\infty(\mathbb R^3\times(0,\infty);\mathbb R^3)}, with support
contained in \NSi{NS-F-P120-008}{\mathcal K\times[0,2]}.
\end{lemma}

\begin{proof}
We choose \NSi{NS-F-P120-009}{\chi_0\in C_c^\infty(\mathbb R)}, equal to one near zero and zero
for arguments at least one. For \NSi{NS-F-P120-010}{\sigma=t-1\geq0}, we set
\NSFormula{NS-F-P120-011}{display}{10.11}
\begin{equation}
f(x,1+\sigma)=\sum_{j=0}^\infty\chi_0(b_j\sigma)\frac{\sigma^j}{j!}F_j(x),
\tag{10.11}\label{eq:10.11}
\end{equation}
where the increasing integers \NSi{NS-F-P120-012}{b_j\geq1} are chosen below. The product rule
and \NSi{NS-F-P120-013}{0\leq\sigma\leq b_j^{-1}} on the support of each summand give
\NSFormula{NS-F-P120-014}{display}{10.12}
\begin{equation}
\left\|\partial_x^\alpha\partial_\sigma^m
 \left(\chi_0(b_j\sigma)\frac{\sigma^j}{j!}F_j\right)\right\|_\infty
\leq C_{j,m}b_j^{m-j}\|F_j\|_{C^{|\alpha|}}.
\tag{10.12}\label{eq:10.12}
\end{equation}
Indeed, when \NSi{NS-F-P120-015}{a} derivatives fall on the cutoff, their factor
\NSi{NS-F-P120-016}{b_j^a} is multiplied by a monomial of degree
\NSi{NS-F-P120-017}{j-m+a}; its bound is
\NSi{NS-F-P120-018}{b_j^{-(j-m+a)}}. Terms whose monomial has been differentiated to zero are
absent. We set \NSi{NS-F-P120-019}{b_0=1}, and for \NSi{NS-F-P120-020}{j\geq1} define
\NSFormula{NS-F-P120-021}{display}{none}
\[
b_j=\min\left\{b\in\mathbb N:b>b_{j-1},\quad
\max_{\substack{a,m\geq0\\a+m\leq\lfloor j/2\rfloor}}
C_{j,m}b^{m-j}\|F_j\|_{C^a}\leq2^{-j}\right\},
\]
where the maximum is over nonnegative integers \NSi{NS-F-P120-022}{a,m}. This is a finite set of
constraints, and \NSi{NS-F-P120-023}{m-j<0} in each of them; thus the defining set of integers is
nonempty. With this choice, \eqref{eq:10.12} is at most
\NSi{NS-F-P120-024}{2^{-j}} whenever
\NSi{NS-F-P120-025}{|\alpha|+m\leq\lfloor j/2\rfloor}.

Every fixed differentiated tail now converges uniformly. At \NSi{NS-F-P120-026}{\sigma=0} the
cutoff in each fixed summand is constant near zero, so its \NSi{NS-F-P120-027}{m}-th derivative
contributes only when \NSi{NS-F-P120-028}{j=m}, giving \NSi{NS-F-P120-029}{F_m}. Uniform
differentiated convergence and \eqref{eq:10.10} show that \eqref{eq:10.11} matches all mixed
space--time derivative limits from \NSi{NS-F-P120-030}{t<1}. Each summand retains support in
\NSi{NS-F-P120-031}{\mathcal K} and is zero for \NSi{NS-F-P120-032}{\sigma\geq1}. The same uniform
convergence proves smoothness at those support boundaries. Finally, the original force is zero near
\NSi{NS-F-P120-033}{t=0}, by \eqref{eq:10.2} and the construction of
\NSi{NS-F-P120-034}{p}.

To state the decay bounds, we choose \NSi{NS-F-P120-035}{R} with
\NSi{NS-F-P120-036}{\mathcal K\subset B(0,R)} and write
\NSi{NS-F-P120-037}{M_{\alpha,m}=\|\partial_x^\alpha\partial_t^mf\|_\infty}. Compact support
gives, for every integer \NSi{NS-F-P120-038}{k\geq0},
\NSFormula{NS-F-P120-039}{display}{none}
\[
|\partial_x^\alpha\partial_t^mf(x,t)|
\leq M_{\alpha,m}(3+R)^k(1+|x|+t)^{-k}.
\]
\end{proof}
